\section{Step 1 Identify the message being hashed}\label{sha-s01}

Begin with the exact three bytes 61, 62 and 63. The displayed word abc does not tell us whether a file also contains a newline. Record the input length and identify where the original message ends before reasoning about padding or digest output.

\subsection{Observe and compare}
Use the completed entries available at this point; keep their inputs and observation boundary explicit. The surrounding explanation identifies the state being compared.

\subsection{Expected observation}
The message length is three bytes; the marker and length field will be added by finalization.

\Needspace{9\baselineskip}
\subsection{Check your understanding}
Is a trailing newline included automatically in this message?

\begin{itemize}
  \item \textbf{A} No; only the declared input bytes belong to it.
  \item \textbf{B} Yes, every text message ends with one.
  \item \textbf{C} Only when the digest begins with a high bit.
\end{itemize}
